\section{Trabalhos Relacionados}
\label{sec:related}

Para a elaboração deste estudo, foram revisados trabalhos comparativos sobre bancos de dados de
séries temporais, abrangendo tanto análises de funcionalidade quanto benchmarks de desempenho.
A busca foi conduzida utilizando palavras-chave como ``data streaming'', ``timeseries'' e
``timeseries databases'', com os resultados ordenados por contagem de citações para priorizar
análises mais estabelecidas e desenvolvidas. A partir dessa busca, 23 resultados foram
selecionados para leitura inicial, a fim de verificar sua relevância para os objetivos deste estudo.

Dentre os 23 trabalhos, alguns não apresentaram relação direta suficiente para uma análise
aprofundada, embora tenham fornecido contexto relevante. \cite{zubaroglu2021data}, por exemplo,
apresenta uma lista de conjuntos de dados aplicáveis a análises de bancos de dados de séries
temporais. \cite{Gaber2005Mining} contribuiu com informações relevantes sobre os algoritmos
subjacentes às implementações de BDSTs, apesar de não estar diretamente focado em avaliação
comparativa.

Da leitura inicial, seis trabalhos foram selecionados para estudo detalhado. Este capítulo está
organizado em três seções. A seção de \textit{Análise de Funcionalidade} cobre trabalhos que
avaliam características técnicas e operacionais dos sistemas BDST. A seção de \textit{Análise
de Desempenho} cobre estudos comparativos de benchmark. A seção de \textit{Análise Comparativa}
sintetiza e critica todos os trabalhos selecionados por meio de uma tabela comparativa,
destacando padrões, lacunas e tendências identificadas na literatura.

\subsection{Análise de Funcionalidade}

\cite{Petre2019} apresentam uma análise baseada em modelo de maturidade de seis bancos de dados
de séries temporais de código aberto: InfluxDB, Graphite, RRDTool, Prometheus, OpenTSDB e
TimescaleDB. O trabalho avalia 18 atributos quantitativos e qualitativos divididos em requisitos
funcionais (aparentes ao usuário final, como capacidades específicas) e não funcionais (não
  diretamente aparentes, mas impactantes, como estratégias de compressão e mecanismos de otimização
de consultas). Os atributos avaliados incluem desenvolvimento ativo, velocidade de desenvolvimento,
maturidade de mercado, replicação de dados e suporte a bibliotecas. O objetivo é auxiliar na
seleção do BDST mais adequado para necessidades de negócio, considerando tanto aspectos técnicos
quanto operacionais. O resultado identificou o InfluxDB como a solução líder na análise.

\subsection{Análise de Desempenho}

Cinco trabalhos no formato de benchmarks controlados ou comparações em cenários reais foram
selecionados para análise detalhada.

\cite{lima2024avaliando} empregaram testes de inserção (variando de 1.000 a 5 milhões de
registros) e consultas de agregação com dados sintéticos (medições indexadas por tempo,
localização, temperatura e umidade) para comparar o desempenho de SGBDs utilizando o Apache
JMeter. A mesma estrutura de dados foi replicada em todos os bancos de dados. Os testes foram
realizados em um ambiente baseado em Docker em uma máquina com AMD Ryzen 5 4500 (6 núcleos),
16\,GB de RAM e armazenamento SSD, testando volumes de 500.000 e 5 milhões de registros. A
reprodutibilidade foi parcial, pois contêineres Docker foram utilizados, mas com detalhes
incompletos para replicação total. O InfluxDB foi o destaque do estudo.

\cite{shah2022performance} compararam o desempenho do InfluxDB~v2.1, TimescaleDB~v2.6,
Druid~v0.22.1 e Cassandra~v3.11.12 em operações de inserção e oito tipos de consulta
(consultas pontuais, agregações e outras), avaliando tempo de carga, consumo de espaço e
latência. Foram utilizados tanto conjuntos de dados reais (dados IoT, dados financeiros, NYC
Taxi) quanto sintéticos (gerados por um gerador Python de séries temporais). O ambiente de
teste consistia em um Intel Core i5-9500 (6 núcleos), 8\,GB de RAM, 500\,GB SSD e
Ubuntu~20.04. O InfluxDB se destacou não apenas no desempenho de consultas, mas também no
consumo de espaço, ficando em segundo lugar apenas após o Druid. O estudo foi classificado
como não reprodutível, pois nenhum script ou configuração foi compartilhado.

\cite{naqvi2017influxdb} compararam InfluxDB, SQL Server, Cassandra, Elasticsearch e OpenTSDB
em três métricas primárias: taxa de ingestão de dados, requisitos de armazenamento e tempo
médio de resposta a consultas. Foram utilizados tanto dados sintéticos (simulando métricas de
servidor) quanto um conjunto de dados real de viagens de táxi em NYC (2009--2017), processados
com ferramentas nativas do SGBD e scripts Java personalizados. Os testes foram realizados em
um Intel Core i5-6300HQ (4 núcleos), 16\,GB de RAM e SSD, usando versões específicas de cada
banco de dados (InfluxDB~1.3, Elasticsearch~5.6.3). O estudo concluiu que o InfluxDB é ideal
para cenários de ingestão de dados temporais de alta velocidade, como IoT e monitoramento, mas
pode não ser adequado para aplicações que exigem operações complexas de JOIN ou atualizações
frequentes. A reprodutibilidade foi classificada como inexistente, sem scripts ou configurações
compartilhadas.

\cite{pereira2023mulletbench} avaliaram o desempenho do InfluxDB~2.7 e Apache IoTDB~1.1.1 em
três cenários: somente nuvem, 2~borda~+~nuvem e 4~borda~+~nuvem, testando cargas de trabalho
de inserção, consulta e mistas, cada uma executada três vezes após reinicialização do sistema.
As ferramentas utilizadas incluíram MulletBench para o benchmark, Ansible e Docker para
orquestração, PCP para monitoramento e tcconfig para simulação de rede. Os testes foram
conduzidos em nós servidores com Intel~i5-9500 (16\,GB de RAM, NVMe SSD) e nós clientes com
Intel~i3-4170 (16\,GB de RAM, SATA SSD) em uma rede comutada Gigabit. Este é o único trabalho
na literatura revisada classificado como totalmente reprodutível, com todas as configurações e
código publicamente disponíveis no GitHub. O Apache IoTDB foi o destaque, superando o InfluxDB
em todos os cenários testados.

\cite{wang2023iotdb} compararam IoTDB~v0.13.0 contra InfluxDB, TimescaleDB, KairosDB, Parquet
e ORC em operações de escrita e leitura, custo de armazenamento e tipos de consulta incluindo
consultas de intervalo, agregações e subamostragem, medindo vazão, latência e tempo de
resposta. Foram utilizados tanto conjuntos de dados industriais reais (dados de veículos, trens
e máquinas) quanto conjuntos sintéticos gerados pelo TSBS. Os testes foram conduzidos em
ambientes standalone (Intel i7-10700, 32\,GB de RAM, HDD de 10\,TB) e distribuídos (Alibaba
Cloud com Xeon Platinum, 256\,GB de RAM), com fator de replicação 3 e particionamento temporal
na configuração distribuída. Embora os detalhes de configuração sejam descritos no artigo,
nenhum guia completo de reprodução foi fornecido, levando a uma classificação de
reprodutibilidade parcial. O Apache IoTDB demonstrou superioridade tanto em ambientes locais
quanto distribuídos, especialmente em grandes escalas industriais envolvendo trilhões de pontos
de dados.

\subsection{Análise Comparativa}

A Tabela~\ref{tab:related} agrega os principais atributos dos cinco estudos de desempenho,
permitindo uma análise comparativa nas dimensões mais relevantes para a metodologia de
benchmark.

\begin{table*}[t]
  \centering
  \caption{Comparação de estudos de benchmark de BDSTs relacionados.}
  \label{tab:related}
  \begin{tabular}{lllllll}
    \hline
    \textbf{Autor} & \textbf{Ano} & \textbf{Bancos de Dados} & \textbf{Reprod.} & \textbf{Datasets} & \textbf{Ferramentas} & \textbf{Líder} \\
    \hline
    \cite{lima2024avaliando}       & 2024 & SQL + NoSQL & Parcial & Sintético          & Apache JMeter                 & InfluxDB \\
    \cite{pereira2023mulletbench}  & 2023 & NoSQL       & Total   & Sintético          & MulletBench, Ansible, Docker  & IoTDB    \\
    \cite{wang2023iotdb}           & 2023 & SQL + NoSQL & Parcial & Real + Sintético   & TSBS                          & IoTDB    \\
    \cite{shah2022performance}     & 2022 & SQL + NoSQL & Nenhuma & Real + Sintético   & timeseries-generator (Python) & InfluxDB \\
    \cite{naqvi2017influxdb}       & 2017 & SQL + NoSQL & Nenhuma & Real + Sintético   & Telegraf, Kapacitor            & InfluxDB \\
    \hline
  \end{tabular}
\end{table*}

Os anos de publicação variam de 2017 a 2024. A maioria dos estudos é recente, com densidade
particular em 2023, refletindo o crescente interesse na avaliação de BDSTs. A presença de
estudos de 2017 destaca que esta não é uma preocupação emergente; no entanto, a rápida evolução
dos próprios bancos de dados significa que descobertas mais antigas podem não se aplicar mais às
versões atuais.

\textbf{Cobertura de bancos de dados.} O InfluxDB está presente em todos os estudos revisados,
servindo como o ponto de referência de facto para comparações de BDSTs, mesmo em trabalhos
onde outro sistema é o assunto principal, como em~\cite{wang2023iotdb}. A presença do
TimescaleDB em três dos cinco estudos reflete sua crescente relevância como alternativa
compatível com SQL. A tendência mais notável é a emergência do Apache IoTDB: aparecendo apenas
nos dois estudos mais recentes (2023), ele lidera em ambos~\citep{pereira2023mulletbench,wang2023iotdb},
sugerindo melhorias de desempenho nas versões recentes, viés de seleção dos benchmarks
mais recentes para cargas de trabalho IoT, ou um crescente reconhecimento de seus pontos
fortes para esse domínio; os dados disponíveis não permitem distinguir entre essas
explicações.

\textbf{Reprodutibilidade.} Esta é a lacuna metodológica mais significativa na literatura
revisada. Dos cinco estudos de desempenho, apenas um (\cite{pereira2023mulletbench})
é totalmente reprodutível, com todo o código e configurações publicamente disponíveis no GitHub.
Dois estudos (\cite{lima2024avaliando} e \cite{wang2023iotdb}) oferecem reprodutibilidade
parcial por meio de contêineres Docker ou configurações descritas, enquanto os dois restantes
(\cite{shah2022performance} e \cite{naqvi2017influxdb}) não fornecem nenhum meio de replicação.
Isso limita a validação independente dos resultados e dificulta atribuir diferenças de
desempenho à arquitetura do banco de dados em vez de às condições experimentais.

\textbf{Conjuntos de dados.} Todos os cinco estudos utilizam conjuntos de dados sintéticos,
refletindo a necessidade de dados controlados e escaláveis em volume. Três estudos incluem
adicionalmente conjuntos de dados reais, como dados de táxi em NYC, leituras de sensores IoT e
dados industriais de veículos~\citep{shah2022performance,naqvi2017influxdb,wang2023iotdb},
proporcionando relevância prática ao lado do controle experimental. O TSBS (Time Series
Benchmark Suite) aparece nos estudos mais recentes, sugerindo uma tendência potencial de
padronização na comunidade.

\textbf{Configurações de hardware.} Os ambientes testados variam consideravelmente: os
processadores variam de Intel~i3 a Intel~i7, a RAM de 8\,GB a 32\,GB (mais nuvem com 256\,GB
em um caso~\citep{wang2023iotdb}), e o armazenamento de SATA SSD a NVMe e HDD. A maioria dos
testes utilizou máquinas de médio porte, consistente com foco em cenários de implantação
realistas. A inclusão de testes em nuvem em estudos mais recentes~\citep{pereira2023mulletbench,wang2023iotdb}
reflete o crescente interesse em arquiteturas borda-nuvem.

\textbf{Tendências de resultados.} O InfluxDB lidera em três dos cinco estudos, todos de
2017--2022. O Apache IoTDB lidera em ambos os estudos de 2023. Esse padrão temporal é notável:
pode indicar melhorias de desempenho nas versões recentes do IoTDB, ou pode refletir o fato de
que benchmarks mais recentes incluem cargas de trabalho específicas de IoT nas quais o modelo
de dados nativo do IoTDB oferece uma vantagem estrutural. A mudança não invalida necessariamente
o InfluxDB para outras cargas de trabalho; no entanto, motiva a inclusão de ambos os bancos de
dados em qualquer avaliação comparativa contemporânea.

Em resumo, a literatura revisada revela três lacunas principais que o presente trabalho aborda:
(1) reprodutibilidade insuficiente: apenas um dos cinco estudos é totalmente replicável;
(2) cobertura limitada do IoTDB junto ao InfluxDB e TimescaleDB no mesmo benchmark controlado;
e (3) a ausência de uma análise sistemática de como as escolhas metodológicas (isolamento de
recursos, ajuste de banco de dados) afetam as conclusões reportadas. Este estudo contribui para
fechar todas as três lacunas.
